% BLOCO P2. Descritivo.
\begin{table}[htbp]\centering
\caption{Sample reconciliation: birth-cohort total vs.\ duration-analysis sample}
\label{tab:P2rec}
\footnotesize
\setlength{\tabcolsep}{4pt}
\begin{tabular}{lrrrr}
\toprule
Cohort & Total & Duration & Excl.$>$47m & \% excl. \\
\midrule
2010 & 128,352 & 106,300 & 20,322 & 15.8 \\
2011 & 146,693 & 126,766 & 18,556 & 12.6 \\
2012 & 155,292 & 139,440 & 14,638 & 9.4 \\
2013 & 153,033 & 140,811 & 11,178 & 7.3 \\
2014 & 151,772 & 146,665 & 4,144 & 2.7 \\
2015 & 146,293 & 146,232 & 0 & 0.0 \\
2016 & 128,143 & 128,143 & 0 & 0.0 \\
\bottomrule
\end{tabular}
\\[3pt]{\footnotesize \textit{Note:} Total: children of the cohort with a valid first application. Duration: those with non-negative eligible time, the sample of the competing-risks analysis. Excl.$>$47m: applicants who applied after the 47-month eligibility limit (age at application above 47 months), who have no eligible time and are dropped; this is the bulk of the gap. The excluded share falls sharply across cohorts because early applicants were much older, so the exclusion is mechanically tied to the timing shift; it is conservative for the trend (adding these children as immediate aging-out lowers early-cohort incidence and widens the gain, see Table~\ref{tab:P2rob}, +Late appl.).}
\end{table}
